% ============================================================================
\chapter{Outillage~: Maven, Gradle, JDK multi-versions, CI}
\label{ch:java-build}
% ============================================================================

\begin{objectif}
\begin{itemize}
  \item Lire et \'ecrire un \fichier{pom.xml} et un \fichier{build.gradle.kts} sans les copier \`a l'aveugle.
  \item Comprendre le cycle de vie, la r\'esolution de d\'ependances, les p\'erim\`etres et les BOM.
  \item Utiliser le wrapper~: \verb|./mvnw| / \verb|./gradlew|--- le build fonctionne chez le voisin.
  \item Produire un ex\'ecutable~: jar, uber-jar, image native GraalVM, Docker.
  \item Formater, analyser et publier~: Spotless, Checkstyle, Error Prone, Release, CI GitHub Actions.
\end{itemize}
\end{objectif}

\section{Cr\'eer le projet}

\begin{terminal}
# Maven : archetypes (le plus simple pour un projet nu)
mvn archetype:generate -DgroupId=com.exemple -DartifactId=app \
    -DarchetypeArtifactId=maven-archetype-quickstart -DarchetypeVersion=1.4 -DinteractiveMode=false
# Gradle (apr\`es installation du CLI)
gradle init --type java-application --dsl kts --test-framework junit-jupiter
# Generateurs en ligne (le plus courant en pratique) :
#   Spring Initializr https://start.spring.io  |  https://start.quarkus.io  |  https://micronaut.io/launch
# Mise a jour des wrappers (a commiter !) :
mvn wrapper:wrapper -Dmaven=3.9.9
gradle wrapper --gradle-version 8.11.1
\end{terminal}

\subsection{La structure~: un contrat, pas une mode}

\begin{codeplain}
app/
\boxv \boxh \boxh mvnw, mvnw.cmd, .mvn/wrapper/          # le wrapper (toujours)
\boxv \boxh \boxh pom.xml            ou   build.gradle.kts + settings.gradle.kts
\boxv \boxh \boxh src/main/java/com/exemple/app/         # code
\boxv \boxh \boxh src/main/resources/                    # config, i18n, schemas SQL, templates
\boxv \boxh \boxh src/test/java/...                       # tests (miroir des paquets)
\boxv \boxh \boxh src/test/resources/                      # jeux de donnees de test
\boxv \boxh \boxh .editorconfig  .gitignore  README.md
\boxe \boxh \boxh (Gradle) build/    (Maven) target/       # JAMAIS versionnes
\end{codeplain}

\section{Maven~: le \texttt{pom.xml} utile}

\begin{codeconf}
<?xml version="1.0" encoding="UTF-8"?>
<project xmlns="http://maven.apache.org/POM/4.0.0"
         xmlns:xsi="http://www.w3.org/2001/XMLSchema-instance"
         xsi:schemaLocation="http://maven.apache.org/POM/4.0.0 https://maven.apache.org/xsd/maven-4.0.0.xsd">
  <modelVersion>4.0.0</modelVersion>

  <!-- 1. Identite : groupId = espace de noms, artifactId = nom, version = semver -->
  <groupId>com.exemple</groupId>
  <artifactId>blog-app</artifactId>
  <version>1.4.0-SNAPSHOT</version>          <!-- -SNAPSHOT = en developpement -->
  <packaging>jar</packaging>

  <!-- 2. Proprietes : encoding, version du langage, versions centralisees -->
  <properties>
    <project.build.sourceEncoding>UTF-8</project.build.sourceEncoding>
    <project.reporting.outputEncoding>UTF-8</project.reporting.outputEncoding>
    <maven.compiler.release>21</maven.compiler.release>   <!-- mieux que source/target -->
    <junit.version>5.11.4</junit.version>
  </properties>

  <!-- 3. BOM : une dependance "import" qui aligne les versions de toute une famille -->
  <dependencyManagement>
    <dependencies>
      <dependency>                            <!-- ne la mets pas dans <dependencies> -->
        <groupId>org.junit</groupId>
        <artifactId>junit-bom</artifactId>
        <version>${junit.version}</version>
        <type>pom</type>
        <scope>import</scope>
      </dependency>
      <dependency>
        <groupId>com.fasterxml.jackson</groupId>
        <artifactId>jackson-bom</artifactId>
        <version>2.18.2</version>
        <type>pom</type><scope>import</scope>
      </dependency>
    </dependencies>
  </dependencyManagement>

  <!-- 4. Depandances -->
  <dependencies>
    <dependency>
      <groupId>org.slf4j</groupId><artifactId>slf4j-simple</artifactId><version>2.0.16</version>
    </dependency>
    <dependency>
      <groupId>org.postgresql</groupId><artifactId>postgresql</artifactId><version>42.7.4</version>
      <scope>runtime</scope>                  <!-- necessaire a l'execution seulement -->
    </dependency>
    <dependency>
      <groupId>org.assertj</groupId><artifactId>assertj-core</artifactId><version>3.27.2</version>
      <scope>test</scope>
    </dependency>
    <dependency>
      <groupId>org.projectlombok</groupId><artifactId>lombok</artifactId><version>1.18.36</version>
      <scope>provided</scope>                 <!-- vu a la compilation, pas embarque -->
    </dependency>
    <dependency>                              <!-- exclusion : l'artefact "le plus faux" -->
      <groupId>commons-logging</groupId><artifactId>commons-logging</artifactId>
      <version>1.3.4</version>
      <exclusions><exclusion><groupId>log4j</groupId><artifactId>log4j</artifactId></exclusion></exclusions>
    </dependency>
  </dependencies>

  <!-- 5. Plugins : le cycle de vie est fait de goals branches sur des phases -->
  <build>
    <finalName>${project.artifactId}</finalName>
    <plugins>
      <plugin>
        <groupId>org.apache.maven.plugins</groupId><artifactId>maven-surefire-plugin</artifactId>
        <version>3.5.2</version>
        <configuration><argLine>-Xshare:off -XX:+EnableDynamicAgentLoading</argLine></configuration>
      </plugin>
      <plugin>
        <groupId>org.apache.maven.plugins</groupId><artifactId>maven-failsafe-plugin</artifactId>
        <version>3.5.2</version>
        <executions><execution><goals><goal>integration-test</goal><goal>verify</goal></goals></execution></executions>
      </plugin>
      <plugin>
        <groupId>org.apache.maven.plugins</groupId><artifactId>maven-shade-plugin</artifactId>
        <version>3.6.0</version>                <!-- uber-jar executable -->
        <executions><execution><phase>package</phase><goals><goal>shade</goal></goals>
          <configuration><transformers>
            <transformer implementation="org.apache.maven.plugins.shade.resource.ManifestResourceTransformer">
              <mainClass>com.exemple.app.App</mainClass>
            </transformer>
          </transformers></configuration>
        </execution></executions>
      </plugin>
      <plugin>                                 <!-- execution : mvn exec:java -->
        <groupId>org.codehaus.mojo</groupId><artifactId>exec-maven-plugin</artifactId><version>3.5.0</version>
        <configuration><mainClass>com.exemple.app.App</mainClass></configuration>
      </plugin>
      <plugin>                                 <!-- versionner proprement -->
        <groupId>org.apache.maven.plugins</groupId><artifactId>maven-release-plugin</artifactId><version>3.1.1</version>
      </plugin>
    </plugins>
  </build>

  <!-- 6. Profils : <profile> avec <id>dev/prod</id>, active par -Pprod ou par activation -->
</project>
\end{codeconf}

\begin{terminal}
mvn clean                # supprime target/
mvn compile              # src/main/java -> target/classes
mvn test                 # tests unitaires (surefire) : src/test/java
mvn package              # jar dans target/
mvn verify               # + tests d'integration (failsafe) + analyses
mvn install              # dans ~/.m2/repository (consomme par les autres projets)
mvn site                 # javadoc, rapports de qualite
mvn -Pprod -DskipTests package
# Diagnostic (les commandes qui te sauvent) :
mvn -X package                                   # log de debug
mvn dependency:tree -Dverbose -Dincludes=com.fasterxml.jackson*   # qui tire quoi
mvn dependency:analyze                            # inutilisees / undeclared
mvn versions:display-dependency-updates ; mvn versions:use-latest-releases
mvn help:effective-pom                            # le POM reel apres heritage et BOM
mvn enforcer:enforce                              # regles (version de Java, interdits)
\end{terminal}

\begin{tablez}[p{0.14\textwidth}p{0.3\textwidth}X]
  \headrow \thd{Scope} \& \thd{Pr\'esent o\`u ?} \& \thd{Effet sur ton artefact} \\
  \midrule
  \verb|compile| (d\'efaut) \& compile + tests + ex\'ecution \& il se retrouve dans le classpath de tes consommateurs : chaque \verb|compile| est une dette publique \\
  \verb|provided| \& compile + tests seulement \& le conteneur ou le JDK le fournit (servlet-api, annotation processors) \\
  \verb|runtime| \& tests + ex\'ecution \& parfait pour un driver JDBC ou une impl\'ementation de log \\
  \verb|test| \& tests seulement \& JUnit, Mockito, Testcontainers, assertions \\
  \verb|system| \& --- \& \`a fuir (jar pos\'e sur le disque) \\
\end{tablez}

\section{Gradle~: le build en Kotlin DSL}

\begin{codeconf}
// build.gradle.kts
plugins {
    java
    application
    `maven-publish`
    id("com.diffplug.spotless") version "6.25.0"
    id("org.graalvm.buildtools.native") version "0.10.4"
    checkstyle
}

group = "com.exemple"
version = "1.4.0"

java {
    toolchain { languageVersion = JavaLanguageVersion.of(21) }   // Gradle telecharge le JDK s'il manque
    withSourcesJar()
}

repositories {
    mavenCentral()
    mavenLocal()                       // pour tes artefacts SNAPSHOT locaux
}

dependencies {
    implementation(platform("org.junit:junit-bom:5.11.4"))
    implementation(platform("com.fasterxml.jackson:jackson-bom:2.18.2"))
    implementation("org.slf4j:slf4j-simple")
    implementation("org.postgresql:postgresql:42.7.4")
    compileOnly("org.projectlombok:lombok:1.18.36")
    annotationProcessor("org.projectlombok:lombok:1.18.36")
    testImplementation("org.junit.jupiter:junit-jupiter")
    testImplementation("org.assertj:assertj-core:3.27.2")
    testRuntimeOnly("org.junit.platform:junit-platform-launcher")
    testImplementation("org.testcontainers:postgresql")
}

application { mainClass.set("com.exemple.app.App") }

testing {
    suites {
        val test by getting(JvmTestSuite::class) { useJUnitJupiter() }
        val integrationTest by registering(JvmTestSuite::class) {
            dependencies { implementation(project()) ; implementation("org.testcontainers:postgresql") }
        }
    }
}

tasks.withType<Test>().configureEach {
    useJUnitPlatform()
    maxParallelForks = (Runtime.getRuntime().availableProcessors() / 2).coerceAtLeast(1)
    testLogging { events("failed", "skipped") }
    reports.html.required = false
}

tasks.withType<JavaCompile>().configureEach { options.encoding = "UTF-8"; options.compilerArgs += "-parameters" }

spotless { java { target("src/*/java/**/*.java"); googleJavaFormat("1.24.0"); removeUnusedImports() } }
check.dependsOn("spotlessCheck")
\end{codeconf}

\begin{codeconf}
// settings.gradle.kts -- multi-module + catalogue de versions centralise
rootProject.name = "blog"
include("app", "domaine", "infra", "api")
dependencyResolutionManagement {
    versionCatalogs {
        create("libs") {
            from(files("gradle/libs.versions.toml"))
        }
    }
}
\end{codeconf}

\begin{codeconf}
# gradle/libs.versions.toml -- le "catalogue" : une seule source de verite
[versions]
junit = "5.11.4"
jackson = "2.18.2"
[libraries]
junit-bom = { module = "org.junit:junit-bom", version.ref = "junit" }
jackson-databind = { module = "com.fasterxml.jackson.core:jackson-databind", version.ref = "jackson" }
[plugins]
spotless = { id = "com.diffplug.spotless", version = "6.25.0" }
\end{codeconf}

\begin{terminal}
./gradlew build                # compile + testes + assemble + verifs (le reflexe)
./gradlew run --args="--port 9090"
./gradlew test --tests "*DepotArticles*" --rerun-tasks
./gradlew --scan               # rapport de build en ligne (Gradle Enterprise/Develocity public)
./gradlew build --offline      # sans reseau (cache deja peuple)
./gradlew build --stacktrace --info
./gradlew dependencies --configuration runtimeClasspath | head -60
./gradlew buildEnvironment ; ./gradlew properties | grep version
./gradlew clean build -x test -PskipChecks
gradle wrapper --gradle-version 8.11.1 --distribution-type bin   # et commite gradle/wrapper/*
# Config globale : ~/.gradle/gradle.properties
org.gradle.jvmargs=-Xmx2g -XX:MaxMetaspaceSize=512m
org.gradle.parallel=true ; org.gradle.caching=true ; org.gradle.configuration-cache=true
org.gradle.daemon=true ; org.gradle.console=rich ; org.gradle.workers.max=8
\end{terminal}

\subsection{Maven ou Gradle~?}

\begin{tablez}[p{0.2\textwidth}p{0.28\textwidth}X]
  \headrow Crit\`ere \& Maven \& Gradle \\
  \midrule
  Courbe d'apprentissage \& douce, convention forte, XML verbeux mais lisible \& plus exigeante (DSL, t\^aches), mais expressive \\
  Vitesse \& correcte, r\'esolutions s\'equentielles \& incr\'ementale + cache + parall\`ele~: 2--10\,\$\times\$ sur les gros builds \\
  Multi-module \& \verb|<modules>| simple \& excellent, graphes de t\^aches fins \\
  Plugins \& \'enorme d\'ep\^ot, versions fig\'ees par super-POM \& \'ecosyst\`eme riche, API de plugin en JVM/Kotlin \\
  Reprisibilit\'e \& \'elev\'ee (XML d\'eclaratif) \& \'elev\'ee \emph{si} le script est disciplin\'e \\
  IDE \& IntelliJ natif, VS~Code correct \& IntelliJ/Android natifs, import parfois capricieux \\
\end{tablez}

\begin{note}
Choisis Gradle pour~: un monorepo, Android, des builds lents, des besoins sur-mesure (g\'en\'eration de
code, Docker, graal). Choisis Maven pour~: un projet serveur classique, une \'equipe h\'et\'erog\`ene, une
entreprise qui centralise d\'ej\`a un d\'ep\^ot Nexus. Les deux se valent techniquement~; la pire option est
le \og ~on a mis les deux~ \fg{}.
\end{note}

\section{Produire un ex\'ecutable~: les quatre formats}

\begin{terminal}
# 1. jar simple (librairie) : mvn package -> target/app.jar ; java -cp app.jar:lib/* com.exemple.App
# 2. uber-jar (tout embarque) -- Maven shade/assembly, Gradle shadow ; attention aux META-INF/services
java -jar app.jar
# 3. distribution avec launcher + classpath (recommande pour un serveur) :
mvn -Pdist package         # ou Gradle : installDist -> build/install/app/bin/app
# 4. image native GraalVM (d\'emarrage ~10 ms, ~40 Mo de RSS, mais build long et r\'eflexion \`a configurer)
nativeCompile          # Gradle : ./gradlew nativeCompile
mvn -Pnative native:compile
native-image -jar app.jar -H:Name=app -march=native --initialize-at-build-time
# 5. Conteneur (la vraie cible en prod) : un Dockerfile multi-etapes
\end{terminal}

\begin{codeconf}
# Dockerfile -- Maven (adapte pour Gradle avec gradle build --no-daemon nativeCompile)
FROM maven:3.9-eclipse-temurin-21 AS build
WORKDIR /m
COPY . .
RUN mvn -B -DskipTests package
FROM eclipse-temurin:21-jre-jammy
RUN useradd -r -u 10001 app && mkdir -p /data && chown app /data
USER app
WORKDIR /app
COPY --from=build /m/target/app.jar .
EXPOSE 8080
HEALTHCHECK --interval=10s --timeout=3s CMD ["sh","-c","wget -qO- http://localhost:8080/health || exit 1"]
ENTRYPOINT ["java","-XX:MaxRAMPercentage=75","-XX:+ExitOnOutOfMemoryError","-jar","/app/app.jar"]
\end{codeconf}

\begin{codeconf}
# .mvn/jvm.config (options de build) ; .mvn/maven.config : -T 1C -Dstyle.color=always
# settings.xml (~/.m2/settings.xml) : miroir, credentials, profils
<settings>
  <mirrors><mirror>
    <id>central-miroir</id><mirrorOf>central</mirrorOf>
    <url>https://maven.aliyun.com/repository/central</url>   <!-- ou un Nexus interne -->
  </mirror></mirrors>
  <servers><server><id>repo-interne</id><username>${env.NEXUS_USER}</username>
           <password>${env.NEXUS_PASS}</password></server></servers>
</settings>
\end{codeconf}

\section{Qualit\'e du build~: la cha\^ine qui rend le code propre tout seul}

\begin{terminal}
# Formatage (le m\^eme partout, sans discussion en revue)
./gradlew spotlessApply            # ou : mvn spotless:apply
mvn com.diffplug.spotless:spotless-maven-plugin:apply
# Analyse statique
mvn -Panalysis verify              # checkstyle + spotbugs + pmd + dependency-check
./gradlew checkstyleMain spotbugsMain pmdMain
mvn org.owasp:dependency-check-maven:check -DfailBuildOnCVSS=7
# Reproduire et debugger le build
mvn -Dmaven.repo.local=/tmp/m2 -o package      # depot local isole / hors ligne
\end{terminal}

\begin{codeconf}
<!-- Checkstyle + PMD + SpotBugs + JaCoCo : les reglages qui rendent le build utile -->
<plugin>
  <groupId>org.apache.maven.plugins</groupId><artifactId>maven-checkstyle-plugin</artifactId><version>3.6.0</version>
  <configuration>
    <configLocation>google_checks.xml</configLocation>       <!-- ou sun_checks.xml, ou le tien -->
    <consoleOutput>true</consoleOutput><failOnViolation>true</failOnViolation>
    <violationSeverity>warning</violationSeverity>
  </configuration>
</plugin>
<plugin>
  <groupId>com.github.spotbugs</groupId><artifactId>spotbugs-maven-plugin</artifactId><version>4.8.3.1</version>
  <configuration><effort>Max</effort><threshold>Medium</threshold><xmlOutput>true</xmlOutput></configuration>
</plugin>
<plugin>
  <groupId>org.jacoco</groupId><artifactId>jacoco-maven-plugin</artifactId><version>0.8.12</version>
  <executions>
    <execution><goals><goal>prepare-agent</goal></goals></execution>
    <execution><id>rapport</id><phase>test</phase><goals><goal>report</goal></goals></execution>
    <execution><id>seuil</id><phase>verify</phase><goals><goal>check</goal></goals>
      <configuration><rules><rule><limits><limit><counter>LINE</counter><minimum>0.80</minimum></limit></limits></rule></configuration>
    </execution>
  </executions>
</plugin>
\end{codeconf}

\section{CI GitHub Actions~: le mod\`ele Java}

\begin{codeyml}
name: build
on:
  push: { branches: [main] }
  pull_request:
jobs:
  java:
    runs-on: ubuntu-latest
    steps:
      - uses: actions/checkout@v4
      - uses: actions/setup-java@v4
        with:
          distribution: temurin
          java-version: 21
          cache: maven            # ou: gradle  (le cache fait gagner 2 a 4 min)
      - name: Build + tests + verification
        run: mvn -B -ntp verify                # ou: ./gradlew build --no-daemon
      - name: Rapport de couverture (JaCoCo)
        run: mvn -B -ntp jacoco:report site -DskipTests
        if: always()
      - uses: actions/upload-artifact@v4
        if: always()
        with:
          name: rapports-${{ matrix.os ?? 'linux' }}
          path: |
            target/site
            target/surefire-reports
            build/reports
  matrice:
    runs-on: ${{ matrix.os }}
    strategy:
      fail-fast: false
      matrix:
        os: [ubuntu-latest, macos-latest, windows-latest]
        java: [17, 21]
    steps:
      - uses: actions/checkout@v4
      - uses: actions/setup-java@v4
        with: { distribution: temurin, java-version: "${{ matrix.java }}" }
      - run: mvn -B -ntp verify
\end{codeyml}

\etape{Publication}
\begin{itemize}
  \item En interne~: \verb|mvn deploy| / \verb|publish| vers un Nexus, Artifactory, GitHub Packages,
        \verb|packagecloud|--- ou un d\'ep\^ot \verb|file:///home/deploy/m2| si tu es seul.
  \item Sur Maven Central~: compte \link{Sonatype Central Portal}{https://central.sonatype.com/},
        groupes valid\'es, \verb|gpg|-signature, javadoc + sources obligatoires,
        \link{central-publishing-maven-plugin}{https://github.com/central-publishing/central-publishing-maven-plugin}
        ou \verb|./gradlew publishToMavenLocal| puis le portal web.
  \item Versionnement~: \verb|x.y.z-SNAPSHOT| en dev~; \verb|mvn release:prepare release:perform|
        (tag \verb|v1.4.0|, passage en \verb|1.4.1-SNAPSHOT|)~; ou \link{Revapi}{https://github.com/revapi/revapi}
        pour refuser une rupture de compatibilit\'e binaire dans une \verb|minor|.
\end{itemize}

\begin{piege}
Cinq erreurs de build classiques~: (1) le wrapper non committ\'e~: chacun a une version de Maven et
les builds diff\`erent~; (2) \verb|maven.compiler.source| sans \verb|release|~: tu compiles avec le
JDK~21 en visant~17 et tu utilises des API de~21~; (3) \verb|<scope>compile| partout~: tes
d\'ependances deviennent publiques et cassent les consommateurs~; (4) deux versions de Jackson
dans l'arbre~: \verb|NoClassDefFoundError| al\'eatoire~; (5) le cache CI sans \verb|~/.m2/repository|.
\end{piege}

\begin{exercice}
\begin{enumerate}
  \item Prends le projet du chapitre pr\'ec\'edent (ou un carnet d'adresses maison) et fais-le passer
        de \og ~\verb|javac| dans un terminal~ \fg{} \`a un projet Maven complet~: BOM, tests, javadoc,
        surefire, failsafe, uber-jar \verb|package|, README avec captures.
  \item Ajoute la m\^eme chose c\^ot\'e Gradle~: toolchain, catalogue de versions, t\^ache
        \verb|int\'egrationTest|~; garde les deux builds un moment et tranche avec des chiffres
        (temps de \verb|clean build| \`a froid et \`a chaud).
  \item Fais \'echouer un build volontairement et r\'epare-le~: d\'ependance en double avec versions
        incompatibles, \verb|NoSuchMethodError|~: puis \verb|dependency:tree|~; un \verb|pom.xml|
        sans UTF-8~: les accents du fichier de propri\'et\'es deviennent illisibles.
  \item Ajoute la cha\^ine qualit\'e~: Spotless (google-java-format), Checkstyle avec une config
        maison de 10 r\`egles, JaCoCo avec seuil bloquant \`a 80 \%, OWASP dependency-check.
        Fais en sorte que \verb|mvn verify| soit ton r\'eflexe local avant chaque push.
  \item \'Ecris le \verb|Dockerfile| multi-\'etapes~: puis une image \verb|-distroless| ou
        \verb|jre-alpine|~: compare taille et temps de d\'emarrage~; ajoute un \verb|HEALTHCHECK|
        et fais tourner le conteneur avec \verb|-XX:MaxRAMPercentage|.
  \item Compile en natif (GraalVM)~: mesure le temps de d\'emarrage et l'empreinte m\'emoire face au
        jar~; documente les r\'eflexions n\'ecessaires (\verb|reflect-config.json|) et ce que tu as
        perdu (diagnostic \`a chaud, JIT qui s'\'echauffe, \verb|jcmd|).
\end{enumerate}
\end{exercice}

\begin{ressources}
\begin{itemize}
  \item Maven --- \url{https://maven.apache.org/guides/} --- Reference Model (le POM, par c\oe{}ur en 2 h)
  \item Gradle --- \url{https://docs.gradle.org/current/userguide/userguide.html} --- Kotlin DSL --- \url{https://docs.gradle.org/current/userguide/kotlin\\\_dsl.html}
  \item \link{Maven Central Search}{https://central.sonatype.com/search} (trouver l'\verb|artifactId| exact)
  \item \link{mvnrepository}{https://mvnrepository.com/} (utile, mais la source de v\'erit\'e reste Central)
  \item GraalVM Native Image --- \url{https://www.graalvm.org/latest/reference-manual/native-image/}
  \item \link{Setup Java action}{https://github.com/actions/setup-java} --- \link{Spotless}{https://github.com/diffplug/spotless}
\end{itemize}
\end{ressources}

\index{Maven}\index{Gradle}\index{wrapper}\index{uber-jar}\index{GraalVM}\index{CI Java}
